\section{Axiomatic Systems and Theory of Proofs}

All of this is our Section 1.3, and that's called axiomatic systems and theory of proofs. You all met proofs---in your physics lectures, theoretical physics lectures, mainly in mathematics lectures, and so on---but what actually is a proof? Well, it's a way of arguing why something is true given something else that you assumed is true. But how do you argue? What are the valid rules of writing down a proof? Well, this can actually be defined, and there are some subtleties in this that actually render some of the proofs that are sometimes given really actually non-proofs. And if you are very serious about proving something, you need to know what the definition of a proof is.

Well, we start with the definition of an axiomatic system. And you see, all this is preparation to write down the axiomatic system that defines set theory.

\textbf{Definition.} An \emph{axiomatic system} is a finite sequence of propositions $a_1, a_2, \ldots, a_n$. Now if one wants to be very smart about it, one could object and say: what do you mean by finite, and what do you mean by numbers $1, 2, 3, 4, 5, \ldots, n$? Because if you haven't defined sets yet, how on earth are you talking about numbers? Well, you could write this down as $a_1 / a_2 / \ldots$---so logicians call these pre-mathematical numbers. It's what you're doing in Bavarian pubs.

So again, if these lectures are attended by a mathematical logician, he will probably shout all the time: well, hang on a second, this is a little bit too na\"ive. It's true---I'm not using the full first-order predicate logic approach to all this thing. It already showed up in my use of tables and so on. But it gives you the right ideas. It's conceptually exactly what they do, without going into too much of the formal nitty-gritty, which for what we're going to do and what we're going to use it for would be fully redundant. Okay, so this is supposed to give you the right idea. It's the right idea; it's not necessarily what a professional mathematical logician would write down if he publishes his new research monograph. Okay, but this is putting the whole thing to classroom use. So this is my disclaimer. Doesn't mean that anything of this is wrong; it's all right.

Okay, so an axiomatic system is a finite sequence of propositions, and these propositions are called the \emph{axioms}.

\textbf{Definition.} A \emph{proof} of a proposition $p$---that's the thing we want to prove---a proof of a proposition within an axiomatic system: so the idea is, I first give you axioms, and then I set a homework problem---prove $p$ starting from the axioms. A proof of proposition $p$ within an axiomatic system $a_1, \ldots, a_n$ is a finite sequence of propositions (or propositional schemes) $q_1, \ldots, q_M$---the steps of the proof, if you wish---such that:
\begin{itemize}
\item $q_M = p$ (the final proposition is supposed to be the thing you want to prove---obviously, if you have a proof and you have steps, at the end you better have the statement that you wanted to prove);
\item for any $j$ that lies between $1$ and $M$, either of the following is true:
\begin{itemize}
\item[\textbf{(A)}] $q_j$ is a proposition from the list of axioms (it means at an arbitrary step in your proof you may pick one of the axioms and put it there---very clear, you may put it there);
\item[\textbf{(T)}] $q_j$ is a tautology;
\item[\textbf{(M)}] there exist $m$ and $n$ somewhere in the proof, but actually before the current step (so $m, n < j$), such that
\[
(q_m \land q_n) \Rightarrow q_j
\]
is true.
\end{itemize}
\end{itemize}

Now we're very careful here. The point is: one single proposition in an axiomatic system can come in the form ``for all $x$, this and that is true,'' and then you can actually use the individual statements. If the proposition is $\forall\, x : P(x)$, it's just a proposition, but you can actually take $P(x)$ out and say: well, it's for all $x$, then I can use $P(x)$ as an individual object, which is then true in maybe infinitely many instances---maybe for all the $x$'s that could be. And that we call a \emph{propositional scheme}. So there is some subtlety there. But the actual steps of the proof must be a finite number of steps. It's very important.

And when I was in my first semester, I had---well, today I would say I was fortunate---that my analysis professor was an intuitionist. So he didn't have proofs by contradiction; that was one thing. And the other thing is that he was sometimes really mad at some textbook proofs of analysis theorems, and he very often said: ``and this proof is not finite.'' And we didn't understand what he meant. Well, we now understand, because in the theory of proofs, you need a finite sequence of propositions. It's very important.

Anyway, so it's such a sequence of steps, and they need to satisfy certain conditions. So, \textbf{(A)} stands for ``axiomatic,'' \textbf{(T)} stands for ``tautology,'' and \textbf{(M)} stands for \emph{modus ponens}, which is a very old logical deduction rule. But we'll write it down here---I mean, there's nothing in the name. Modus ponens: there exist $m$ and $n$ before the current step, so if we're in the $j$-th step of the proof we may use any of the previous steps---step $m$ and step $n$, which are before the current step---such that $(q_m \land q_n) \Rightarrow q_j$ is true.

So please realize: this is a proposition. Okay, so it could be that we have two steps before in the proof, and if I put them together by ``and'' and I can show that from this, via the implication arrow, $q_j$ is true---for some $m$ and some $n$ before this $j$-th step, this whole thing is true---that also qualifies the $j$-th step as a valid one. So the $j$-th step can consist of an axiom, it can consist of a plain tautology, or it can be deduced from two previously arrived-at steps.

I think I didn't define what a tautology is. A \emph{tautology} is a proposition that is always true, independent of the elementary propositions it's made up of. E.g., $p \lor \lnot p$. This is a proposition, right? And because $p$ appears in there, and we don't know whether $p$ is true or false, generically we wouldn't know whether $p \lor \lnot p$ is true or false---because it could depend, would depend normally, on $p$. But this specific combination $p \lor \lnot p$---well, this is always true, as you can check from the table. This is called a tautology.

Yes---do you really need the last one, or doesn't it follow out of A and T somehow? I mean, it's pretty much just the assembly of propositions, this\ldots\ Okay, let's take no. You need it. You definitely, seriously, desperately, fully need it. It's just---probably somewhere---it's just---sorry to intrude---you don't really see anything, yes, because you have some kind of kindergarten intuition about this thing. You see, if this is true, then I can do this. But remember, this is a binary operator. It's not more. It has not this nice intuition about where the arrow goes, or the thought goes, or something. It's just a binary operator. This is a logical statement, as one of the axioms will be.

At the moment we don't have enough technology to provide a meaningful proof. But in this course I will write out one proof in full according to this scheme---one single one---and that is, I think, the uniqueness of the empty set. Now there will be an axiom that stipulates there exists an empty set with such-and-such properties, and then you have to show there's only one such set. It's a very quick proof, and intuitively we say ``yeah, it's clear''---that, if you know the definition of the implication arrow. But if you really want to write it out, I think it has something like 16 lines. It's a 16-step proof if you write it out in full like this. And I want to present one proof that is precise, and the other proofs will be given in the form they're usually given, by some quick arguments.

But the understanding of mathematicians is---I mean, even in the best journals, the most celebrated theorems are usually not proven in this fashion, because this goes into the millions and billions, I don't know, of lines, proving comparatively simple stuff. Okay, so we don't do that. This is kind of like machine code in computer science, right? Almost like the zeros and the ones, the binary code. And actually the proofs we give are in the higher language, like C++, also. Okay, so this is what we do in mathematics. But the understanding is that the proofs we write in books---we say ``well, now it's obvious that'' blah blah blah---the understanding is it's only a proof if your fellow mathematicians agree that in principle what was written down now could probably, according to our intuition, be broken down to proofs of this type. Okay?

And if somebody has the feeling that what you do here will not lead to a proof of this type---for instance, not to a finite-step proof, or a proof with finitely many steps---then mathematicians will raise objection to this proof. They will say: well, we're not convinced. And then in principle, as long as somebody says ``I'm not convinced,'' a mathematician has to break down the proof more and more, until you arrive here. So you can obstruct any lecture in mathematics---so don't tell anyone---you can obstruct any lecture in mathematics by\ldots\ you know.

Okay---no, it doesn't matter, because you can go back in your proof to an arbitrary point, because the rough idea is that everything there has been there anyway: it's an axiom, or it's a tautology, and all the other stuff has been arrived at by implication arrows. That's the intuition. And so it suffices to look at stuff that was before. But because everything that comes after has been connected by an implication arrow, you cannot go further than the current step, because that would imply that you used equivalence arrows. If you use equivalence arrows, you can go further down---say ``oh, and later I show, so I can use it again.'' Now apart from some circularity, it's the directedness of this thing, very roughly speaking, intuitively speaking, that allows you to go to any previous pair of previous steps, even if there's something in between, but not further. Okay, fine.

\textbf{Remark.} If $p$ can be proven from an axiomatic system $a_1, \ldots, a_n$, we often write:
\[
a_1, \ldots, a_n \vdash p,
\]
where $\vdash$ means ``proves.'' That's a short form for this.

And of course an obvious remark is that this definition of proof allows you to \emph{recognize} a proof. In principle, you can give this to a computer. You say: computer, please check, is the following a proof of this? That's very easy to check. An altogether different matter is to \emph{find} a proof for a certain proposition. And as you all know from your own experience---even set problems where usually, if you're asked to prove something, it can be proven---it's not so easy. And sometimes it's so difficult that if you find a proof where people before didn't, you become a celebrated mathematician. So the theory of proofs here doesn't give you any hint of how to prove things better, but at least it tells you: what you have written down right now, that is a proof. And that can be very useful, of course.

A further remark is that obviously any tautology, should it occur in an axiomatic system---if among the axioms there's a tautology---can be removed from the list of axioms without impairing the power of the axiomatic system. And it's very clear: it means you can still prove the same stuff, because if you don't pull your tautology from the list of axioms, you can pull it anyway---so it's actually this axiom \textbf{(T)} that guarantees that.

And an extreme case of this is---that's why I'm mentioning it---what is the axiomatic system for propositional logic? What is the list of axioms there? It's the \emph{empty sequence}. Because in propositional logic, all we can prove are tautologies. Want to show something is true? Well, but then you don't need any axioms, because you can always pull a tautology into your proof by axiom \textbf{(T)}.

But let's have one further definition.

\textbf{Definition.} An axiomatic system is called \emph{consistent} if there exists a proposition $Q$ which cannot be proven from the axiomatic system. Formally:
\[
\text{not } (a_1, \ldots, a_n \vdash Q).
\]
There exists such a $Q$.

Now what's the idea behind this surprising, maybe, definition? Well, it's very simple. Consider an axiomatic system containing contradictory propositions. So for instance, $a_1, \ldots$, whatever, at some point we have $S$ as one of the propositions, and later on in the list we have $\lnot S$. So we have an axiomatic system that stipulates both $S$ and $\lnot S$.

Then by the deduction rule \textbf{(M)}---the one whose existence or justification was doubted tentatively---by \textbf{(M)}, clearly: I can pull $S$ into my proof, I can pull $\lnot S$ into my proof by axiom \textbf{(A)}, and then I can use these two by \textbf{(M)} and connect them with an ``and.'' I only need to check whether, for an arbitrary $Q$,
\[
(S \land \lnot S) \Rightarrow Q
\]
is a tautology. Yeah, this is a tautology, because $S \land \lnot S$ is always false. You remember: but if the assumption is always false, the implication arrow is defined such that the whole statement is always true. \emph{Ex falso quodlibet.} So it's again the definition of the implication arrow, which I dwelled on at the beginning so much, which shows that this is true. \emph{Ex falso quodlibet}---whatever you like, whatever statement you like, can be concluded. So by \textbf{(M)} we can arrive at the truth of any statement $Q$; we can also arrive at the statement $\lnot Q$. It doesn't matter. So the problem is: any statement can be proven if you have contradictory axioms.

Now it's a sign of not having inconsistency---of not having contradictory axioms---if it's simply not true that every statement can be proven. But here, every statement can be proven. Hence we need to exclude that every statement can be proven. This is just one example; the above is the general definition of consistency.

Well, having come this far, we can now write down an impressively sounding theorem, which is very trivial.

\textbf{Theorem.} Propositional logic is consistent.

That's reassuring. What's the proof? Well, we have to---it suffices to show that there exists a proposition that cannot be proved within propositional logic. That's clear: there's at least one that cannot be proven, then it's consistent.

Now propositional logic has an empty sequence of axioms. You don't have any. So the only rules for proof---only \textbf{(T)} and \textbf{(M)}---must carry any proof. Now let's look at \textbf{(T)} and \textbf{(M)} again. The only way by which you can introduce an entirely new step is by writing down a tautology, and then the modus ponens only allows you to take previous tautologies and to write down the $q_j$ where all this thing is a tautology. So the only thing you can prove are tautologies. Only tautologies can be proven. But that means that $Q \land \lnot Q$ cannot be proven, because it is not a tautology. So we constructed the proposition that cannot be proven, and so propositional logic is consistent. \qed

And this is really trivial. Now again, the important remark is that while it's perfectly fine and clear how to define consistency, it's perfectly difficult to prove it for a given axiomatic system---propositional logic being a big exception. But if the axiomatic system becomes more powerful, if you have more axioms---like we're going to write down for set theory---you might then again ask the question: is this set of axioms consistent in this sense? And it's very difficult to prove that.

In fact, there are famously these statements, first arrived at by G\"odel, that say, to a certain extent, it's even not possible to show this under certain circumstances. And I will give only a very rough taste of one of the things G\"odel proved.

\textbf{Theorem} (G\"odel). Roughly---not being technically precise---it goes as follows. Any axiomatic system that is powerful enough to encode the elementary arithmetic of the natural numbers (which are the positive integers)---so you see, I use imprecise language here. What's elementary arithmetic? Well, in elementary arithmetic, it depends on what you need to prove: it means you understand prime numbers, you understand products, stuff like that. Okay, rather elementary stuff---let's say elementary school mathematics. What exactly it is, one would see in the proof, what you need. But any axiomatic system that's powerful enough to encode elementary arithmetic of natural numbers is either inconsistent---never mind its successes with the arithmetic of the natural numbers---is either inconsistent, or contains a proposition that can neither be proven nor disproven.

So variants of this formulation are: it contains true statements that cannot be proven. I mean, this sent shockwaves through the mathematics world at the time. This was the beginning of the 20th century---1920s, 1930s, something like that---because in mathematics, of all the arts if you wish, the notion of truth seemed to be clear up to then. Truth---or, something was true if it could be proven based on a system of assumptions. This was the pure truth of mathematics. But then G\"odel came along and he ironically proved this, and variants of this.

Now how does the proof proceed? Well, the proof is involved, but the basic idea of the proof is the following. First, assign to each mathematical, or also meta-mathematical, statement or proposition a number---the so-called, what later on people called, \emph{G\"odel number}. Well, a mathematical statement is ``$a^2 + b^2 = c^2$,'' something of this kind---``for all'' blah blah, ``there exists,'' and so on. And G\"odel, very roughly speaking, says: the plus sign gets a number, the exponent gets a number, every variable gets a different number, the equal sign gets a different number. And then you can start talking about this equation. You can say: the left-hand side of this equation has more symbols than the right-hand side. That would be a meta-mathematical statement, because obviously, ``the left-hand side has more symbols than the right-hand side'' is not a mathematical statement; it's a meta-mathematical statement. But even to such meta-mathematical statements, G\"odel was able---or devised how---to assign a number.

Take primes. Everything---every elementary symbol---is assigned a prime, and then you take powers of primes and you multiply. The powers get big numbers. And then to meta-mathematical statements you get a huge number---boom, there it is, right? Now you do this for just everything. Okay, and then once you did that, in a sense you have a number for everything. And because you have to multiply the numbers, you need elementary arithmetic. Now you understand why there's this mysterious condition on the axiomatic system that you're using---being able to handle the elementary arithmetic---because this assignation of numbers and the calculation and so on requires this elementary arithmetic.

And then what you do: you use a ``the barber shaves all people in his village who do not shave themselves'' type argument to identify a proposition that is neither provable nor disprovable. So you use some reflexive construction. You know this, I think, from the proof that there are uncountably many real numbers. You know that proof: you assume that you can write down every real number in the form integer, decimal point, and then decimal places $d_1, d_2, d_3, d_4, \ldots$, and then you assume that you can number every such expression---you assign a number to them. But then you actually use the number of the expression to construct a new expression, a new real number, which you show is not in the previous list of numbers. It's kind of reflexive. You might want to revisit this proof. This is of a similar reflexive type. And here you do the same thing: you assign to statements G\"odel numbers, and then you produce a statement that uses the G\"odel number of the very same statement. Okay, and then you arrive at such a contradiction.

And as you can see, we defined ``axiomatic system,'' we defined what it is to be ``provable,'' and thereby therefore what it means to be ``disprovable''---that means you can prove the negation---and we defined what ``consistent'' is. So according to our definitions, we fully understand this theorem; we understand this statement. The proof is a different matter.

Okay, so at this place I will rest with the development of the language we're going to use to define axiomatic set theory, which we will start next time. And we will start by looking at the element symbol $\in$, which is nothing but a predicate of two variables, because it takes a variable on the left---$x \in y$---and a variable on the right. It's just a matter of notation to write the symbol in the middle rather than writing $\in(x, y)$. Okay, it's just notational convenience, and we very often do this for predicates of two variables: we write the name of the predicate in the middle, because it just saves you two brackets and a comma.

But this will be the elementary predicate of two variables, and there is no other in set theory. Even the equal sign between sets will be defined in terms of the element sign. And from this one elementary predicate of two variables---which will not be further explained; there will be nobody explaining to you in axiomatic set theory what it actually means that $x$ is an element of $y$---and indeed, $x$ and $y$ will both be sets. So there's no such hierarchy like ``you are a set, a grown-up set, but you're only an element.'' No, no---sets can be contained in sets. And a priori, one could even think about the question: is a set an element of itself? There's no a priori reason why you would exclude this. It would give you some kind of recursion. And you could ask: can you construct a set of all the sets that do not contain themselves? If you don't like these sets containing themselves, you could think: well, let's constrain the set universe to all those sets that do not contain themselves. And then you very quickly---actually in two lines---prove that this construction (we will do this next time) is not a set.

The set---or let me call it the collection---of all the sets that do not contain themselves: you only use this element predicate. This is no longer a set, as you can prove by asking the question whether that set contains itself or not---and that is neither provable nor disprovable; both lead to a contradiction. So that immediately shows: okay, so then set theory is tricky, and we need a system of axioms that regulates our use of the element predicate. Okay, because we know nothing of it; we just know there it is. And now the axioms give you rules of how to deal with it. It's a kind of an operational definition. Okay, you see: if it behaves like this, or if you use it like this, then you get what we want. We can construct sets.

There are some axioms to tell you: these and those sets exist. For instance, there exists an empty set. Or: if you have a set that exists, the power set of the set---that's the set of all the subsets of the set---you have to postulate this by axiom. And all the other sets we're actually going to construct from these previous sets. Well, in fact, you can show that the only set you really need to postulate is the empty set. So in principle, every set we're ever going to talk about is ultimately constructed from the empty set, the set that contains the empty set, the set with the element ``empty set,'' stuff like this.

Okay, but there are actually, I think, nine axioms. We'll see; it depends on how you write them down; the number varies slightly. There are about nine axioms that you can write down that regulate set theory, and that reproduce the set theory on which all of modern mathematics is built. And occasionally, you have an axiom that's called the \emph{axiom of choice}, for instance---famously---or Zorn's lemma, some people call it. The axiom of choice: and you wonder, well, do you really want to have that? And why do we need to be so picky about the axioms? There we go: for instance, if you do not postulate the axiom of choice, you cannot prove that every vector space has a basis. Now this we use very often.

So you see, already in our building down there, where we're talking about sets, already there are some screws you have to justify in order to have the whole building stand as we want it to stand. Okay, and there are other examples for that. So at least for completeness---and that's the purpose of today's lecture and of the next lecture---is to really look at the foundations of the subject, to have a good and justifiably good feeling about the further constructions that are to come.

See you next time, with axioms of set theory. Thank you.
